%% GENERATED by Thesis/tools/make_appendices.py -- do not edit by hand.
%% Every constant below is read from the script that uses it.
\chapter{Complete Parameter Configuration}
\label{app:hyperparameters}

Everything required to reproduce a run appears here, taken from the values the pipeline applies and not transcribed by hand. Where a value differs across configurations the table shows the variation rather than a representative value, because two of those variations are threats to validity discussed in Chapter~\ref{ch:discussion} and hiding them behind an average would be the wrong choice.

\section{Constant Across All 27 Configurations}

These are the values the audit pins. A configuration that changed any of them would not be one of the 27 this dissertation reports, because the comparison rests on their being identical instead of on their being any particular value.

\begin{table}[htbp]
\centering
\small
\caption{Parameters identical in every configuration, enforced by the audit described in Section~\ref{sec:method-overview}.}
\label{app:tab-shared}
\begin{tabularx}{\textwidth}{@{}l l X@{}}
\toprule
Parameter & Value & Where it acts \\
\midrule
\texttt{TRAINING\_FRAC} & 0.7 & chronological train/evaluation split \\
\texttt{THRESHOLD\_PERCENTILE} & 95.0 & label-blind decision threshold \\
Regression parent cap & 3 & parents per target, chosen by shortest lag \\
Attribution depth, tabulated & 3 & variables reported per attack \\
Attribution depth, plotted & 15 & bars per attribution figure \\
RLS warm-up & 15 steps & before drift is scored \\
\texttt{ALPHA} & 0.05 & significance level for discovery \\
\bottomrule
\end{tabularx}
\end{table}

\section{Per-Domain and Per-Algorithm Settings}

The subsampling rate and the variable count are properties of each algorithm's workflow and not of the pipeline, and the asymmetry they produce is the study's primary threat to internal validity. It is shown here in full rather than summarised.

\begin{table}[htbp]
\centering
\small
\caption{Subsampling rate, variable count, forgetting factor and graph sparsification multiplier, by algorithm and domain.}
\label{app:tab-domain}
\begin{tabular}{lllrrrr}
\toprule
Domain & Algorithm & Variables & Subsample & $\lambda$ & $\kappa$ & Sparsify \\
\midrule
Robotics & PCMCI & 244 & 74 & 0.995 & 1.2 & 1.0 \\
 & GES & 35 & 10 & 0.995 & 1.2 & 1.0 \\
 & CAM-UV & 35 & 10 & 0.995 & 1.2 & 1.0 \\
\midrule
Healthcare & PCMCI & 12 & 7 & 1.0 & 2.0 & 1.0 \\
 & GES & 12 & 10 & 1.0 & 2.0 & 1.0 \\
 & CAM-UV & 12 & 10 & 1.0 & 2.0 & 1.0 \\
\midrule
Industrial & PCMCI & 52 & 1 & 1.0 & 3.0 & 0.0 \\
 & GES & 52 & 10 & 1.0 & 3.0 & 0.0 \\
 & CAM-UV & 52 & 10 & 1.0 & 3.0 & 0.0 \\
\bottomrule
\end{tabular}
\end{table}

\section{Ridge Penalty by Domain and Regression Family}

The full grid is given because the penalty varies by regression variant on one domain, which partially confounds RQ2 there. It is identical across the three algorithms within every cell, which is why RQ1 is unaffected; the audit check pinning regression knobs across algorithms enforces that invariant on every build.

\begin{table}[htbp]
\centering
\caption{Ridge penalty $\alpha$ for each domain and regression family, identical across all three algorithms within each cell.}
\label{app:tab-ridge}
\begin{tabular}{lrrr}
\toprule
Domain & Linear / Ridge & Polynomial & RBF + Ridge \\
\midrule
Robotics & 10.0 & 0.01 & 1.0 \\
Healthcare & 10.0 & 10.0 & 10.0 \\
Industrial & 10.0 & 10.0 & 10.0 \\
\bottomrule
\end{tabular}
\end{table}

\section{Algorithm and Feature-Map Settings}

These are the settings that belong to one algorithm or one feature map rather than to the pipeline, so they do not appear in the uniformity tables above.

\begin{table}[htbp]
\centering
\caption{Settings specific to individual algorithms and to the non-linear feature maps.}
\label{app:tab-algo}
\begin{tabular}{lll}
\toprule
Component & Parameter & Value \\
\midrule
PCMCI & independence test & ParCorr \\
PCMCI & $\tau_{\max}$ & derived per domain: 2, 3, 3 \\
GES & score & causal-learn default BIC \\
GES & edge rule & both adjacency entries agree the direction is determined \\
CAM-UV & $\alpha$ & 0.05 \\
CAM-UV & confounder mode & integrate \\
Polynomial & degree & 3 \\
RBF & $\gamma$ & 0.001 \\
RBF & components & 50 \\
\bottomrule
\end{tabular}
\end{table}

\section{Neural Baseline Settings}

The protocol below is applied identically to every architecture on every domain, which is what makes the baseline comparison in Chapter 6 a tuned one, not a default one.

\begin{table}[htbp]
\centering
\caption[Training protocol shared by the three autoencoder baselines]{Training protocol shared by the three autoencoder baselines. The search grid is applied identically to every architecture on every domain.}
\label{app:tab-neural}
\begin{tabular}{ll}
\toprule
Parameter & Value \\
\midrule
Sequence length & 10 \\
Search grid & $h \in \{32, 64\}$, $l \in \{1, 2, 3\}$, lr $\in \{0.001, 0.01\}$ \\
Selection & lowest fault-free validation loss \\
Validation split & 0.15 \\
Maximum epochs & 300 \\
Early-stopping patience & 15 \\
Seeds & 42, 43, 44 \\
Optimiser & Adam \\
\bottomrule
\end{tabular}
\end{table}

\section{Hardware and Software Environment}

The machine is recorded because Table~\ref{tab:runtime} reports wall-clock times measured on it, and those figures mean nothing without it. Every result in this dissertation was produced on an Intel Core i5-8250U at 1.60~GHz with four cores, eight threads and 32~GB of memory, running Windows 11. It is a 15~W mobile processor and not a workstation: a reader reproducing this work on faster hardware should expect the wall-clock figures to fall and the findings, which rest on AUC-ROC, not on time, to be unaffected. The package versions below are those installed in the environment that produced the results reported here.

\begin{table}[htbp]
\centering
\caption{Package versions.}
\label{app:tab-env}
\begin{tabular}{ll}
\toprule
Package & Version \\
\midrule
python & 3.13.3 \\
numpy & 2.5.2 \\
scipy & 1.18.0 \\
scikit-learn & 1.9.0 \\
pandas & 3.0.5 \\
torch & 2.13.0 \\
shap & 0.52.0 \\
matplotlib & 3.11.1 \\
tigramite & 5.2.10.1 \\
causal-learn & 0.1.4.4 \\
lingam & 1.13.0 \\
\bottomrule
\end{tabular}
\end{table}
